\section{引言}
量子色动力学（QCD）描述夸克—胶子在很宽能量范围内的相互作用，并在低能区提出
非微扰挑战。强子结构研究依赖于部分子分布函数（PDF），而 PDF 的确定又依赖于
精确的 QCD 计算，影响着 LHC 以及未来项目（EIC、HL-LHC、LHeC）的研究。
对新物理的探寻取决于精确的实验—理论比较，这凸显了 PDF 精度的重要性。

对撞机过程中的 QCD 因子化，除 PDF 外，还涉及广义部分子分布（GPD）和
横向动量依赖分布（TMD）。虽然这些分布函数不能直接测量，
但可以通过把过程的截面因子化为硬与软两部分来确定它们。

在深度非弹性散射（DIS）中，结构函数 $F_k(x,Q^2)~(k=1,2,3)$ 扮演关键角色，
它们依赖 $Q^2=-q^2$，其中 $q$ 是转移到核子上的 4 动量，
$x = Q^2/2(pq)$ 表示部分子携带的核子动量份额。这里 $p$ 为核子动量。
当 $Q^2$ 远大于核子质量时，因子化定理~\cite{Collins:1989gx}
把结构函数与标度依赖的部分子分布函数 $f_a(x,\mu^2)$
同 Wilson 系数 $C_k^a(Q^2/\mu^2)$ 的卷积联系起来
（对夸克味与胶子 $a$ 求和）。
Wilson 系数可以微扰计算，PDF 随 $\mu$ 的演化由
Dokshitzer-Gribov-Lipatov-Altarelli-Parisi（DGLAP）
方程~\cite{Gribov:1972ri,Lipatov:1974qm,Dokshitzer:1977sg,Altarelli:1977zs} 描述。
只要给定某标度 $\mu=\mu_0$ 处的初条件，DGLAP 方程即可求解。
PDF 的演化可微扰计算，但实验数据的精度总是要求更高阶的计算。
重正化与因子化方案的标准选择是 $\MSbar$ 方案。
通过分析一组涉及核子的硬散射过程，可以从实验上确定 PDF。
长期以来，学界一直致力于用定义明确的 $x$ 依赖参数化来确定 PDF，
并从 $\mu_0^2 = 1 - 4$GeV 附近开始演化。

格点 QCD 可为 PDF 及其演化的确定提供理论输入。
由于欧几里得几何以及运动学中光锥的支配地位，
直接在格点 QCD 中计算 PDF 存在特殊困难。
PDF 与强子矩阵元（可在格点 QCD 中计算）之间的联系
通过 PDF 的矩 $\left\langle x^{n} \right\rangle$ 建立~\cite{Curci:1980uw,Collins:1981uw}。
由文献~\cite{Kronfeld:1984zv,Martinelli:1987zd,Martinelli:1987bh} 开创的
PDF 矩的格点 QCD 计算，原则上提供了完整重构 PDF 的途径。
但由于计算高阶矩所面临的理论与数值困难，
这一可能性一直不切实际。为确定 PDF 的 $x$ 依赖，
已发展出若干替代方法：重夸克 OPE~\cite{Aglietti:1998mz,Detmold:2005gg}、
准部分子分布（quasi-PDF）~\cite{Ji:2013dva}、
赝 PDF（pseudo-PDF）~\cite{Radyushkin:2016hsy,Karpie:2018zaz}、
基于 OPE 的方法~\cite{Chambers:2017dov}、流—流方法~\cite{Braun:2007wv,Ma:2014jla}，
以及强子张量方法~\cite{Liang:2019frk}。
值得注意的是，这些方法原则上也能间接确定核子 PDF 的矩
~\cite{HadStruc:2021qdf,Gao:2022uhg} 与介子的矩~\cite{Gao:2020ito,Gao:2022iex}。

不同的高能散射过程对应不同的 PDF：非极化、螺旋度与横向性（transversity）。
某些运动学区域以及特定的 PDF 难以从实验数据中提取。
这凸显了从格点 QCD 模拟数据重构 PDF 的重要意义。

在本工作中，我们重新审视计算 PDF 矩的可能性；若能成功，
将为从 QCD 确定分布函数提供一条重要且互补的途径。
我们描述的方法同时解决了过去阻碍从格点 QCD 计算任意阶矩的
理论与数值两方面的挑战。
